\documentclass[11pt,a4paper]{article}
\input{preamble}
\title{Asymptotic Methods}
\begin{document}
\maketitle
\thispagestyle{empty}

\noindent\textbf{Pre-requisites}
\begin{itemize}
    \item Complex methods / analysis.
    \item Refer to Section II in Stuart.
\end{itemize}

\section*{1. Intro \& Motivation}

\begin{itemize}
    \item Most real problems don't have closed form solutions, or if they do, they are unworkable or difficult to interpret.

    \item Numerical approximations can be useful, but they can be expensive, they may have bad numerical properties, or solutions may be inscrutable.

    \item Taylor series, etc.\ can be limited in applicability.

    \item Asymptotic methods emphasise usefulness, interpretability as well as accuracy.
\end{itemize}

\noindent\textbf{Examples}

\medskip

\noindent\textbf{e.g. 1)}
\[
\Gamma(z)=\int_0^\infty t^{z-1}e^{-t}\,dt,
\qquad \operatorname{Re}(z)>0.
\]

For $n$ an integer,
\[
\Gamma(n+1)=n!
\sim \sqrt{2\pi}\,n^{n+\frac12}e^{-n}
\left(
1+\frac{1}{12n}+\frac{1}{288n^2}+\cdots
\right)
\qquad \text{as }n\to\infty.
\]

\medskip

\noindent\textbf{e.g. 2)}
\[
x^2y''+xy'-x^2y=0.
\]

Exact series solution:
\[
I_0(x)=\sum_{n=0}^{\infty}
\frac{x^{2n}}{2^{2n}(n!)^2}.
\]

Liouville--Green:
\[
\Longrightarrow
I_0(x)\sim(2\pi x)^{-1/2}e^x
\qquad \text{as }x\to\infty.
\]

\medskip

\noindent\textbf{e.g. 3)}
\[
y''+k^2e^{-x^2/2}y=0,
\]
with
\[
y\to y_0 \qquad \text{as }x\to\infty,
\]
\[
y\to \pm y_0 \qquad \text{as }x\to-\infty.
\]

\[
k=?,
\qquad k=k_0,k_1,k_2,\ldots
\]

\[
\Longrightarrow
k_n\sim\frac{\sqrt{\pi}}{2}
\left(n+\frac12\right)
\qquad \text{as }n\to\infty.
\]

\noindent WKBJ.

\medskip

\begin{itemize}
    \item Convergent methods typically give arbitrarily accurate approximations with more \& more terms/iterations.

    \item Asymptotic methods give arbitrarily accurate approximations as some variable approaches a limit. Thus they always come with a qualifier, e.g.\ $z\to0$ or $n\to\infty$.
\end{itemize}


\section*{2. Asymptotic Expansions}

\subsection*{2.1 Definitions \& Notation}

\subsubsection*{(a) Asymptotic Equality/Equivalence}

For 2 functions $f(x)$ and $g(x)$, we have
\[
f(x)\sim g(x)
\qquad \text{about }x=x_0
\]
if
\[
\boxed{
\lim_{x\to x_0}\frac{f(x)}{g(x)}=1.
}
\]

\noindent\textbf{Remarks:}

\begin{itemize}
    \item Usually we want to approximate $f$ with $g$.

    \item We can write
    \[
    f(x)=
    \underbrace{g(x)}_{\text{dominant}}
    +
    \underbrace{[f(x)-g(x)]}_{\text{subdominant}}.
    \]

    \item Note $\sim$ is an equivalence relation
    (that is reflexive, transitive, symmetric).

    \item It respects products:
    \[
    f_1\sim g_1,\qquad f_2\sim g_2
    \]
    \[
    \Longrightarrow f_1f_2\sim g_1g_2.
    \]

    \item It is linear:
    \[
    \alpha f_1+\beta f_2
    \sim \alpha g_1+\beta g_2,
    \]
    where $\alpha,\beta$ are constants.

    \item Definition has problems when dealing with oscillatory functions that go through $0$ as $x\to x_0$.
\end{itemize}


\subsubsection*{(b) ``Little $o$''}

Suppose $f$ \& $g$ are real functions.

We write
\[
\boxed{
f(x)=o(g(x))
\qquad \text{as }x\to x_0
}
\]
if
\[
\boxed{
\lim_{x\to x_0}\frac{f(x)}{g(x)}=0.
}
\]

\noindent\textbf{Remark:}

\[
f(x)=o(g(x))
\]
can also be written as
\[
f(x)\ll g(x)
\qquad \text{as }x\to x_0.
\]

\begin{itemize}
    \item Can be framed in inequalities, e.g.\ $x_0=\infty$ and let $g>0$.

    Then $f=o(g)$ as $x\to\infty$ is equivalent to
    \[
    \forall\varepsilon>0,\quad
    \exists a
    \quad\text{s.t. when }x>a,
    \]
    \[
    \boxed{|f|\leq\varepsilon g.}
    \]
\end{itemize}


\subsubsection*{(c) ``Big $O$''}

We write
\[
f(x)=O(g(x))
\qquad \text{as }x\to x_0
\]
for $g>0$, if
\[
\exists C>0
\quad\text{s.t.}
\]
\[
\boxed{
|f(x)|\leq Cg(x)
}
\]
for $x$ sufficiently close to $x_0$.

\medskip

\noindent\textbf{Warning:}

In some areas of physics people use big-$O$ and $\sim$ interchangeably.


\subsubsection*{Little Exercises: Show That}

\noindent\textbf{(a)}
\[
f(x)\sim g(x)
\iff
f-g=o(g)
\qquad \text{as }x\to x_0.
\]

\medskip

\noindent\textbf{(b)}
\[
\sin x\ll x^2
\qquad \text{as }x\to\infty.
\]

\medskip

\noindent\textbf{(c)}
\[
12x^5+6x^3+x=O(x^5)
\qquad \text{as }x\to\infty.
\]

\medskip

\noindent\textbf{(d)}
\[
e^{ax}\ll e^{bx}
\qquad \text{as }x\to\infty,
\]
if $b>a>0$.

\medskip

\noindent\textbf{(e)}
\[
e^x\gg x^N
\qquad \text{as }x\to\infty,
\]
\[
\forall N\in\mathbb{Z}^+.
\]


\subsection*{2.2 Asymptotic Expansions}

\begin{itemize}
    \item Have unworkable/unknown $f(x)$.

    \item Want to know its behaviour in a limit.

    \item Want to find a simpler/known $g(x)$ so that
    \[
    f\sim g.
    \]

    \item Also want to break up $g$ into a hierarchy of approximations, s.t.
    \[
    g(x)=\sum_{k=0}^{N}a_k\phi_k(x),
    \]
    where $a_k$ are coefficients, $\phi_k$ are known \& workable, and where adding more terms gives a better approximation that can be estimated.
\end{itemize}


\subsubsection*{(a) Asymptotic Sequence}

A sequence of functions
\[
\{\phi_0(x),\phi_1(x),\ldots\}
\]
is an ``asymptotic sequence'' as $x\to x_0$ if
\[
\boxed{
\phi_{k+1}(x)\ll\phi_k(x)
\qquad \text{as }x\to x_0,
\quad \forall k.
}
\]

\medskip

\noindent\textbf{e.g. 1)}
\[
\phi_k=(x-a)^k
\qquad \text{around }x=a.
\]

Check:
\[
\frac{\phi_{k+1}}{\phi_k}
=x-a\to0
\qquad \text{as }x\to a.
\quad \checkmark
\]

\medskip

\noindent\textbf{e.g. 2)}
\[
\phi_k=e^{-kx}x^k
\qquad \text{as }x\to\infty
\text{ or }x\to-\infty.
\]

Check:
\[
\frac{\phi_{k+1}}{\phi_k}
=e^{-x}x.
\]

\[
\begin{aligned}
e^{-x}x&\to0
&&\text{as }x\to\infty
&&\checkmark\\
e^{-x}x&\to\infty
&&\text{as }x\to-\infty
&&\times
\end{aligned}
\]

\medskip

\noindent\textbf{e.g. 3)}
\[
\phi_k=\cos(nx)
\qquad \text{as }x\to\infty.
\]

\[
\times
\]

\medskip

\noindent\textbf{e.g. 4)}
\[
\left\{
|x|,\,
\log|x|,\,
\log|\log|x||,\,
\ldots
\right\}
\qquad \text{as }x\to\infty.
\]

\[
?
\]


\subsubsection*{(b) Asymptotic Expansion}

\begin{itemize}
    \item Suppose $f(x)$ is a function we want to understand near $x=x_0$.

    \item Let $\{\phi_k\}_{k=0}^{\infty}$ be an asymptotic sequence.
\end{itemize}

Then, the finite sum
\[
\sum_{k=0}^{N}a_k\phi_k(x)
\]
for some coefficients, is an ($N$-term) asymptotic approximation of $f(x)$ as $x\to x_0$ if:
\[
\boxed{
f(x)=\sum_{k=0}^{N}a_k\phi_k(x)+o(\phi_N(x))
}
\qquad \text{as }x\to x_0.
\]

Alternatively, we can write the above as
\[
\boxed{
\frac{f(x)-\displaystyle\sum_{k=0}^{N}a_k\phi_k(x)}
{\phi_N(x)}
\longrightarrow 0
}
\qquad \text{as }x\to x_0.
\]

\begin{center}
    \textbf{End of lecture 1}
\end{center}

\subsection*{(c) Asymptotic expansions/series}

If (b) holds for \emph{all} \(N\), then this is called an
``asymptotic series or expansion'', and we write
\[
f(x) \sim \sum_{k=0}^{\infty} a_k \phi_k(x)
\qquad \text{as } x \to x_0.
\]

\noindent
\textbf{Attention:} No claim of convergence is made. It is just notation.
We never fix \(x\) and then take \(N \to \infty\).

Instead, we have an infinite number of approximations (of different \(N\)),
and \(x \to x_0\) is taken in each separately!

\subsection*{(d) Properties of asymptotic expansions and remarks}

\subsubsection*{(I) Mathematical operations}

Suppose \(\{\phi_k(x)\}\) is an asymptotic sequence as \(x \to x_0\), and let
\[
f \sim \sum_{k=0}^{\infty} a_k \phi_k,
\qquad
g \sim \sum_{k=0}^{\infty} b_k \phi_k.
\]

\noindent
\textbf{Linearity}
\[
\alpha f + \beta g
\sim \sum_{k=0}^{\infty} (\alpha a_k + \beta b_k)\phi_k
\qquad \text{as } x \to x_0.
\]

\noindent
\textbf{Multiplication}

If \(\phi_k = (x-x_0)^k\), then
\[
fg \sim \sum_{k=0}^{\infty} c_k \phi_k,
\qquad
c_k = \sum_{j=0}^{k} a_j b_{k-j}.
\]

\noindent
\textbf{Division}

See 3.8 BO.

\medskip
\noindent
\textbf{Integration}

Term-by-term integration is usually okay.

For example, if \(\phi_k = (x-x_0)^k\), then
\[
\int_{x_0}^{x} f(t)\,dt
\sim \sum_{k=0}^{\infty}
\frac{a_k(x-x_0)^{k+1}}{k+1}
\qquad \text{as } x \to x_0.
\]

\noindent
\textbf{Exercise:} Use the definitions to prove this.

\medskip
\noindent
\textbf{Differentiation}

Term-by-term differentiation is not always okay.

To do this, you need \(f'(x)\) to be continuous and to have an
asymptotic series, i.e.
\[
f'(x) \sim \sum_{k=0}^{\infty} c_k \phi_k(x).
\]

\subsubsection*{(II) Different asymptotic representations}

Suppose \(\{\phi_k\}\) and \(\{\psi_k\}\) are different asymptotic
sequences as \(x \to x_0\). Also,
\[
f \sim \sum_{k=0}^{\infty} a_k \phi_k,
\qquad
f \sim \sum_{k=0}^{\infty} b_k \psi_k
\qquad \text{as } x \to x_0.
\]
Then \(a_k\) and \(b_k\) are not the same.

For example,
\[
f(x) = 1 + \frac{1}{x}.
\]
Consider \(\phi_k = x^{-k}\) as \(x \to \infty\). Then
\[
f(x) \sim \sum_{k=0}^{\infty} a_k \phi_k,
\qquad
a_0 = 1,\quad a_1 = 1,\quad a_n = 0 \quad (n>1).
\]

Now consider
\[
\psi_k = (1+x)^{-k}.
\]
Then
\[
\begin{aligned}
f
&= 1 + \frac{1}{x}
 = \frac{1}{1-(x+1)^{-1}} \\
&= 1 + \frac{1}{1+x} + \frac{1}{(1+x)^2} + \cdots \\
&\sim \sum_{k=0}^{\infty} b_k \psi_k,
\end{aligned}
\]
where
\[
b_k = 1 \qquad \forall k.
\]

\subsubsection*{(III) Uniqueness of coefficients}

Let \(\{\phi_k\}\) be an asymptotic sequence, and
\[
f(x) \sim \sum_{k=0}^{\infty} a_k \phi_k
\qquad \text{as } x \to x_0.
\]
Then the coefficients \(a_k\) are unique.

In fact, the \(a_k\) can be constructed:
\[
a_0 = \lim_{x \to x_0} \frac{f(x)}{\phi_0(x)},
\]
\[
a_n = \lim_{x \to x_0}
\left[
\frac{f(x)-\displaystyle\sum_{k=0}^{n-1} a_k\phi_k(x)}
{\phi_n(x)}
\right].
\]

\noindent
\textbf{Proof.} First, we know
\[
\lim_{x \to x_0}
\frac{f(x)-\displaystyle\sum_{k=0}^{N} a_k\phi_k}
{\phi_N}
= 0
\qquad \forall N.
\]

Let \(N=0\). Then
\[
\lim_{x \to x_0}
\frac{f(x)-a_0\phi_0}{\phi_0}
= 0,
\]
\[
\lim_{x \to x_0}
\left(\frac{f(x)}{\phi_0}\right)-a_0
= 0.
\]

Use induction on \(N\).

Suppose the formulae are true for \(N=n\), and the coefficients are
unique for \(k \le n\). Then
\[
\begin{aligned}
\lim_{x \to x_0}
\frac{f(x)-\displaystyle\sum_{k=0}^{n} a_k\phi_k}
{\phi_n}
&= 0 \\
&=
\lim_{x \to x_0}
\left(
\frac{f(x)-\displaystyle\sum_{k=0}^{n-1} a_k\phi_k}
{\phi_n}
\right)-a_n.
\end{aligned}
\]

\noindent
Alternative proof: see p.~3 in Manton.

\subsubsection*{(IV) An asymptotic series does not uniquely determine
the function it is describing}

For example, suppose
\[
f(x) \sim g(x)
\qquad \text{as } x \to x_0,
\]
and
\[
h(x) \ll f(x)
\qquad \text{as } x \to x_0.
\]
Then
\[
f(x)+h(x) \sim g(x)
\qquad \text{as } x \to x_0.
\]
So, if \(g\) is an asymptotic expansion, it can describe both
\(f\) and \(f+h\).

For example,
\[
f_1 = \frac{1}{1+x},
\qquad
f_2 = \frac{1+e^{-x}}{1+x}.
\]
Both are
\[
\sim \sum_{k=0}^{\infty} (-1)^k x^{-(k+1)}
\qquad \text{as } x \to \infty.
\]

\subsubsection*{(V) Power series}

In this course, we mostly take
\[
\phi_k = (x-x_0)^k
\qquad \text{as } x \to x_0,
\]
or
\[
\phi_k = x^{-k}
\qquad \text{as } x \to \infty.
\]

In fact, a Taylor series about a point is an asymptotic series about
that point:
\[
f(x) -
\sum_{k=0}^{N}
\frac{f^{(k)}(a)(x-a)^k}{k!}
= o\!\left(|x-a|^N\right).
\]

\noindent
\textbf{Exercise:} Prove this using the definitions and the integral
form of the remainder term.

\medskip
\noindent
\textbf{Terminology:} A term is said to be ``small beyond all orders''
in some limit if it cannot be represented by any power series or,
equivalently, is smaller than any power.

For example, \(e^{-x}\) is small beyond all orders in inverse powers
of \(x\) as \(x \to \infty\):
\[
e^{-x} \ll x^{-k}
\qquad \forall k,
\qquad \text{as } x \to \infty.
\]

\subsubsection*{(VI) Optimal truncation}

\noindent
\textbf{Question:} If an asymptotic series diverges, how many terms
should I include in the sum?

If
\[
f(x) \sim \sum_{n=0}^{\infty} a_n \psi_n(x)
\qquad \text{as } x \to x_0,
\]
and we want to get the best approximation to \(f(x)\) at \(x=x_1\),
where \(x_1\) is close to \(x_0\), then truncate the sum at the
smallest term (or the one before).

So stop adding terms when
\[
|a_k\psi_k(x_1) |<| a_{k+1}\psi_{k+1}(x_1)|.
\]

This gives you an unexpectedly excellent approximation.

\begin{center}
\textbf{End of lecture 2}
\end{center}

\subsection*{2.3 A. expansions in $\mathbb{C}$ --- Stokes phenomenon}

\begin{itemize}
    \item So far our defns relied on real limits, $x,x_0\in\mathbb{R}$ etc.
    \item When we use complex variables, $x=z\in\mathbb{C}$ and $z\to z_0$,
    typically $z$ must be restricted to a sector of $\mathbb{C}$ around $z_0$
    for some of our definitions to work.

    [See III.2 in Stuart]

    e.g.\ little $o$ \& $O$ defns need to bring in moduli of functions.
    \item Even if $f(z)$ analytic at $z_0$, it may have different a.\ series
    as $z_0$ is approached from different sectors in $\mathbb{C}$.
    This is Stokes phenomenon.
\end{itemize}

\begin{center}
    \pdfximage width 0.42\linewidth {figures/asymptotic-lecture3-sectors.png}
    \pdfrefximage\pdflastximage
\end{center}

e.g.
\[
    f(z)=\sinh(z^2)=\frac{1}{2}\left(e^{z^2}-e^{-z^2}\right),
    \qquad z\in\mathbb{C}.
\]
What is its behaviour as $|z|\to\infty$?

Which of $e^{z^2}$ or $e^{-z^2}$ is dominant?
\[
    \left|\frac{e^{-z^2}}{e^{z^2}}\right|
    =\left|e^{-2z^2}\right|=e^{-2\operatorname{Re}(z^2)}.
\]

\noindent
(a) Suppose $|z|\to\infty$ \& $\operatorname{Re}(z^2)>0$.
\[
    \text{ratio}\to0
    \quad\Rightarrow\quad e^{-z^2}\ll e^{z^2}
    \quad\Rightarrow\quad f(z)\sim\frac{1}{2}e^{z^2}
    \quad\text{in this limit.}
\]
\[
    \operatorname{Re}(z^2)>0
    \quad\Rightarrow\quad
    \Omega:\quad
    \begin{gathered}
        -\frac{\pi}{4}+\delta<\arg z<\frac{\pi}{4}-\delta,\\
        \text{and}\quad
        \frac{3\pi}{4}+\delta<\arg z<\frac{5\pi}{4}-\delta,
    \end{gathered}
\]
for small $\delta$.

\begin{center}
    \pdfximage width 0.72\linewidth {figures/asymptotic-lecture3-stokes-lines.png}
    \pdfrefximage\pdflastximage
\end{center}

\noindent
(b) $|z|\to\infty$ \& $\operatorname{Re}(z^2)<0$.
\[
    \text{ratio}\to\infty
    \quad\Rightarrow\quad e^{z^2}\ll e^{-z^2}
    \quad\&\quad f\sim-\frac{1}{2}e^{-z^2}.
\]
\[
    \widetilde{\Omega}:\quad
    \frac{\pi}{4}+\delta<\arg z<\frac{3\pi}{4}-\delta
    \quad\&\quad
    -\frac{3\pi}{4}+\delta<\arg z<-\frac{\pi}{4}-\delta.
\]
\[
    f(z)\sim
    \begin{cases}
        \frac{1}{2}e^{z^2}, & |z|\to\infty\ \&\ z\in\Omega,\\
        -\frac{1}{2}e^{-z^2}, & |z|\to\infty\ \&\ z\in\widetilde{\Omega}.
    \end{cases}
\]

\begin{itemize}
    \item A.\ behaviour of $f$ as $|z|\to\infty$ depends on which wedge
    $z\to\infty$ within.
    \item $\mathbb{C}$ is divided into sectors, separated by special lines,
    called \underline{Stokes lines}.

    [Warning, sometimes called anti-Stokes lines!!]
\end{itemize}

\section*{3. Approximations of Integrals}

\subsection*{3.1 Taylor expanding the integrand --- Stieltjes integrals}

(Warm up problem)

Consider:
\[
    I(x)=\int_0^\infty\frac{e^{-t}}{1+xt}\,dt,
    \qquad x\geq0.
\]
\begin{itemize}
    \item We want to understand $I$'s behaviour when $x\to0^+$.
    \item How about we expand $\dfrac{1}{1+xt}$?
\end{itemize}
\[
    \frac{1}{1+xt}=\sum_{n=0}^\infty(-1)^n x^n t^n
    \qquad\text{\Large$\times$}\quad t\text{ can be big!}
\]

Consider instead
\[
    \sum_{n=0}^N(-xt)^n
    =\frac{1-(-xt)^{N+1}}{1+xt}
    =\frac{1}{1+xt}-\underbrace{\frac{(-xt)^{N+1}}{1+xt}}_{r_N}.
\]
\[
    \Rightarrow\quad\frac{1}{1+xt}=\sum_{n=0}^N(-xt)^n+r_N.
\]
\begin{align*}
    I(x)&=\int_0^\infty e^{-t}
    \left[\sum_{n=0}^N(-1)^n x^n t^n+r_N\right]dt\\
    &=\sum_{n=0}^N(-1)^n x^n
    \underbrace{\int_0^\infty t^n e^{-t}\,dt}_{=n!}+R_N\\
    &=\sum_{n=0}^N(-1)^n n!\,x^n+R_N\\
    &=\sum c_n\phi_n(x)+R_N.
\end{align*}
$\phi_n=x^n$ as $x\to0^+$, so is a.\ sequence!

Need to show $R_N=o(x^N)$ as $x\to0^+$.
\[
    \begin{aligned}
        \frac{R_N}{x^N}
        &=\frac{x^{N+1}}{x^N}\int_0^\infty
        \frac{t^{N+1}e^{-t}}{1+xt}\,dt\\
        &\leq x\int_0^\infty t^{N+1}e^{-t}\,dt
        =x(N+1)!\longrightarrow0
        \quad\text{as }x\to0^+.
    \end{aligned}
\]
Note $1+xt>1$.
True $\forall N$.
\[
    \Rightarrow\quad I(x)\sim\sum_{n=0}^\infty(-1)^n n!\,x^n
    \qquad\text{as }x\to0^+.
\]

\subsection*{3.2 Integration by parts}

Recap:
\[
    \int_a^b f'(t)g(t)\,dt
    =[fg]_a^b-\int_a^b fg'\,dt.
\]
Generalisation:
\begin{align*}
    I&=\int_a^b f(t)g^{(n)}(t)\,dt\\
    &=\sum_{k=1}^n(-1)^{k-1}
    \left[f^{(k-1)}(t)g^{(n-k)}(t)\right]_a^b
    +(-1)^n\int_a^b f^{(n)}(t)g(t)\,dt.
\end{align*}

\begin{itemize}
    \item Repeated integration by parts can generate an a.\ approximation
    under certain circumstances.
\end{itemize}
e.g.\ How does
\[
    F(x)=\int_x^\infty e^{-t^2/2}\,dt
\]
behave as $x\to\infty$?
\begin{align*}
    F(x)&=\int_x^\infty\frac{te^{-t^2/2}}{t}\,dt\\
    &=\left[-\frac{1}{t}e^{-t^2/2}\right]_x^\infty
    +\int_x^\infty\frac{e^{-t^2/2}}{t^2}\,dt\\
    &=\underbrace{\frac{e^{-x^2/2}}{x}}_{\text{a.\ approximation}}
    +\underbrace{\int_x^\infty\frac{e^{-t^2/2}}{t^2}\,dt}_{\text{error}}.
\end{align*}
Need to check that
\[
    \text{error}=o\!\left(\frac{e^{-x^2/2}}{x}\right)
    \qquad\text{as }x\to\infty.
\]
\begin{align*}
    \text{error}&=\int_x^\infty\frac{te^{-t^2/2}}{t^3}\,dt
    \qquad(\text{note }t\geq x)\\
    &\leq\frac{1}{x^3}\int_x^\infty te^{-t^2/2}\,dt
    =\frac{1}{x^3}e^{-x^2/2}\\
    &=o\!\left(\frac{e^{-x^2/2}}{x}\right)
    \qquad\text{as }x\to\infty.
\end{align*}
\[
    \Rightarrow\quad F(x)\sim\frac{e^{-x^2/2}}{x}
    \qquad\text{as }x\to\infty.
\]

In fact we find
\[
    F(x)\sim\frac{e^{-x^2/2}}{x}
    \sum_{n=0}^\infty(-1)^n\frac{(2n-1)!!}{x^{2n}}
    \qquad\text{as }x\to\infty.
\]
\[
    (2n-1)!!=(2n-1)(2n-3)(2n-5)\cdots3\cdot1.
\]
Exercise: Show this.

\begin{center}
    \textbf{End of lecture 3}
\end{center}

\noindent\textbf{Example 2.}
Find the behaviour of
\[
    I(x)=\int_a^b f(t)e^{xt}\,dt
    \qquad\text{as }x\to\infty,
\]
where $b>a$, $f(b)\neq0$, and $f$ is smooth on $[a,b]$.

As $x\to\infty$, the integral is dominated by a narrow region near $t=b$
because of the exponential. Very roughly,
\[
    I(x)\approx f(b)\int_a^b e^{xt}\,dt
    =\frac{f(b)}{x}(e^{xb}-e^{xa})
    \sim\frac{f(b)e^{xb}}{x},
\]
using $e^{xb}\gg e^{xa}$ as $x\to\infty$.

Repeated integration by parts, using
$e^{xt}=\dfrac{d}{dt}\left(\dfrac{e^{xt}}{x}\right)$, gives
\begin{align*}
    I(x)
    &=\sum_{k=1}^N(-1)^{k-1}
      \left[f^{(k-1)}(t)e^{xt}x^{-k}\right]_a^b
      +\frac{(-1)^N}{x^N}\int_a^b f^{(N)}(t)e^{xt}\,dt\\
    &=e^{xb}\bigl(J_1-J_2+J_3\bigr),
\end{align*}
where
\begin{align*}
    J_1&=\sum_{k=1}^N(-1)^{k-1}f^{(k-1)}(b)x^{-k},\\
    J_2&=e^{-x(b-a)}
         \sum_{k=1}^N(-1)^{k-1}f^{(k-1)}(a)x^{-k},\\
    J_3&=\frac{(-1)^N}{x^N}
         \int_a^b f^{(N)}(t)e^{x(t-b)}\,dt.
\end{align*}
$J_1$ could be an $N$-term a.\ approximation with $\phi_k=x^{-k}$
as $x\to\infty$, but we need to show that
\[
    J_2=o(x^{-N}),
    \qquad J_3=o(x^{-N})
    \qquad\text{as }x\to\infty.
\]
For the first term,
\[
    \frac{J_2}{x^{-N}}
    =e^{-x(b-a)}
      \sum_{k=1}^N(-1)^{k-1}f^{(k-1)}(a)x^{N-k}
    \longrightarrow0
    \qquad\text{as }x\to\infty,
\]
which is true for every $N$. For the remainder,
\begin{align*}
    \left|\frac{J_3}{x^{-N}}\right|
    &\leq\int_a^b |f^{(N)}(t)|e^{x(t-b)}\,dt\\
    &\leq\max_{a\leq t\leq b}|f^{(N)}(t)|
          \int_a^b e^{x(t-b)}\,dt\\
    &=\max_{a\leq t\leq b}|f^{(N)}(t)|
      \frac{1-e^{-x(b-a)}}{x}
      \longrightarrow0
      \qquad\text{as }x\to\infty.
\end{align*}
This holds for every $N$ for which $f$ has $N$ continuous derivatives.
If $f$ is infinitely differentiable, then
\[
    I(x)=\int_a^b f(t)e^{xt}\,dt
    \sim\sum_{k=1}^\infty
    (-1)^{k-1}\frac{f^{(k-1)}(b)e^{xb}}{x^k}
    \qquad\text{as }x\to\infty.
\]

\subsection*{3.3 The Gamma function}

Recall
\[
    \Gamma(z)=\int_0^\infty t^{z-1}e^{-t}\,dt,
    \qquad \operatorname{Re}(z)>0.
\]
\begin{enumerate}
    \renewcommand{\labelenumi}{(\alph{enumi})}
    \item $\Gamma$ is analytic for $\operatorname{Re}(z)>0$.
    \item $\Gamma(z+1)=z\Gamma(z)$ \quad (recursion relation).
    \item $\Gamma(1)=1$ and $\Gamma(n+1)=n!$ for positive integers $n$.
    \item For $x>0$ and $\operatorname{Re}(z)>0$,
    \[
        \int_0^\infty t^{z-1}e^{-xt}\,dt
        =\frac{\Gamma(z)}{x^z}.
    \]
    \item Generalisation of binomial coefficients:
    \[
        \binom{z}{k}
        =\frac{1}{k!}\frac{\Gamma(z+1)}{\Gamma(z-k+1)},
        \qquad z\in\mathbb{C},\quad k=0,1,2,3,\ldots,
    \]
    with the quotient interpreted by continuation where necessary.
    \[
        (1+t)^z=\sum_{k=0}^\infty\binom{z}{k}t^k,
        \qquad |t|<1.
    \]
    \item Reflection and duplication formulae: see Section 3.2 in
    Chapling.
\end{enumerate}

\subsection*{3.4 Watson's lemma}

Consider
\[
    F(x)=\int_0^T f(t)e^{-xt}\,dt
    \qquad\text{as }x\to\infty,
    \qquad 0<T\leq\infty.
\]
Suppose
\[
    f(t)\sim t^\alpha\sum_{n=0}^\infty a_n t^{\beta n}
    \qquad\text{as }t\to0^+,
    \qquad \alpha>-1,\quad\beta>0.
\]
Also assume that either
\begin{enumerate}
    \renewcommand{\labelenumi}{(\roman{enumi})}
    \item $|f(t)|<Ke^{bt}$ for all $t\in[0,T]$, for some $K,b>0$; or
    \item $\displaystyle\int_0^T |f(t)|\,dt<\infty$.
\end{enumerate}
Then
\[
    \boxed{
    F(x)\sim\sum_{n=0}^\infty
    a_n\frac{\Gamma(\alpha+\beta n+1)}{x^{\alpha+\beta n+1}}}
    \qquad\text{as }x\to\infty.
\]
This is \emph{Watson's lemma}.

\noindent\textbf{Example 2.}
\[
    xy''+(a+1-x)y'-ay=0,
    \qquad a>0.
\]
One solution is
\[
    M(x)=\frac{e^x}{\Gamma(a)}
    \int_0^1 e^{-xt}(1-t)^{a-1}\,dt.
\]
What is its behaviour as $x\to\infty$?

Here $f(t)=(1-t)^{a-1}$ is integrable, so case (ii) applies, and
\[
    f(t)=(1-t)^{a-1}
    =\sum_{n=0}^\infty
    (-1)^n\frac{\Gamma(a)}{\Gamma(a-n)n!}\,t^n,
    \qquad |t|<1.
\]
Thus $\alpha=0$, $\beta=1$, and Watson's lemma gives
\begin{align*}
    M(x)
    &\sim\frac{e^x}{\Gamma(a)}
    \sum_{n=0}^\infty
    \frac{(-1)^n\Gamma(a)\Gamma(n+1)}{\Gamma(a-n)n!}
    \frac{1}{x^{n+1}}\\
    &=e^x\sum_{n=0}^\infty
    \frac{(-1)^n}{\Gamma(a-n)x^{n+1}}
    \qquad\text{as }x\to\infty.
\end{align*}
In particular,
\[
    M(x)\sim\frac{e^x}{x\Gamma(a)}
    \qquad\text{as }x\to\infty.
\]

\noindent\textbf{Proof of Watson's lemma.}
Write
\begin{align*}
    F(x)
    &=\int_0^\epsilon f(t)e^{-xt}\,dt
      +\int_\epsilon^T f(t)e^{-xt}\,dt\\
    &\equiv F_1(x)+F_2(x),
\end{align*}
where $\epsilon>0$ is small and fixed, with $\epsilon<T$.

First, show that $F_2$ is small.

\noindent Case (i): for $x>b$,
\begin{align*}
    |F_2(x)|
    &\leq\int_\epsilon^T Ke^{bt}e^{-xt}\,dt\\
    &\leq\int_\epsilon^\infty Ke^{-(x-b)t}\,dt\\
    &=\frac{K}{x-b}e^{(b-x)\epsilon}\\
    &=O\!\left(\frac{e^{-\epsilon x}}{x}\right)
    \qquad\text{as }x\to\infty.
\end{align*}

\noindent Case (ii):
\begin{align*}
    |F_2(x)|
    &\leq\int_\epsilon^T |f(t)|e^{-xt}\,dt\\
    &\leq e^{-\epsilon x}\int_\epsilon^T |f(t)|\,dt\\
    &\leq e^{-\epsilon x}\int_0^T |f(t)|\,dt\\
    &=O(e^{-\epsilon x})
    \qquad\text{as }x\to\infty.
\end{align*}
In both cases, $F_2$ is small beyond all orders in inverse powers of $x$.

\begin{center}
    \textbf{End of lecture 4}
\end{center}
\subsection*{3.3 Gamma function (again)}

$\Gamma$ can be analytically continued to $\mathbb{C}$, except at
$z=0,-1,-2,\ldots$, using
\[
    \Gamma(z)=\frac{\Gamma(z+1)}{z}.
\]

\noindent\textbf{Proof of Watson's lemma (continued).}
Recall that
\[
    F(x)=\int_0^\epsilon f(t)e^{-xt}\,dt
        +\int_\epsilon^T f(t)e^{-xt}\,dt
        =F_1(x)+F_2(x),
\]
with $F_2(x)=O(e^{-\epsilon x})$ as $x\to\infty$.
We want to show, for every $N$, that
\[
    F_1(x)=\sum_{n=0}^N
    a_n\frac{\Gamma(\alpha+\beta n+1)}{x^{\alpha+\beta n+1}}
    +o\!\left(x^{-\alpha-\beta N-1}\right).
\]
Using property (d) of the Gamma function, consider
\begin{align*}
    &\left|F_1(x)-\sum_{n=0}^N
    a_n\frac{\Gamma(\alpha+\beta n+1)}{x^{\alpha+\beta n+1}}\right|\\
    &\quad=\left|\int_0^\epsilon f(t)e^{-xt}\,dt
    -\int_0^\infty\sum_{n=0}^N a_n t^{\alpha+\beta n}e^{-xt}\,dt\right|\\
    &\quad\leq\int_0^\epsilon e^{-xt}
    \left|f(t)-\sum_{n=0}^N a_n t^{\alpha+\beta n}\right|\,dt
    +\left|\int_\epsilon^\infty e^{-xt}
    \sum_{n=0}^N a_n t^{\alpha+\beta n}\,dt\right|\\
    &\quad\equiv R_1+R_2.
\end{align*}
For $R_2$, put $u=t-\epsilon$. Then
\begin{align*}
    R_2
    &\leq e^{-\epsilon x}\sum_{n=0}^N |a_n|
    \int_\epsilon^\infty e^{-x(t-\epsilon)}t^{\alpha+\beta n}\,dt\\
    &=e^{-\epsilon x}\sum_{n=0}^N |a_n|
    \int_0^\infty e^{-xu}(\epsilon+u)^{\alpha+\beta n}\,du.
\end{align*}
For each $n$, there exist $C_n,k_n>0$ such that
\[
    (\epsilon+u)^{\alpha+\beta n}<C_n e^{k_nu}
    \qquad\text{for all }u\in[0,\infty).
\]
Hence, for $x>\max_{0\leq n\leq N}k_n$,
\begin{align*}
    R_2
    &\leq e^{-\epsilon x}\sum_{n=0}^N C_n|a_n|
    \int_0^\infty e^{(k_n-x)u}\,du\\
    &=e^{-\epsilon x}\sum_{n=0}^N\frac{C_n|a_n|}{x-k_n}
    =O\!\left(\frac{e^{-\epsilon x}}{x}\right)
    \qquad\text{as }x\to\infty.
\end{align*}

In $R_1$, use the asymptotic expansion of $f$. In particular,
\[
    f(t)-\sum_{n=0}^N a_n t^{\alpha+\beta n}
    =o(t^{\alpha+\beta N})
    =O(t^{\alpha+\beta(N+1)})
    \qquad\text{as }t\to0^+,
\]
where the last bound follows by taking one more term in the expansion.
By the definition of $O$, for sufficiently small $\epsilon$ there is
$K_3>0$ such that
\begin{align*}
    R_1
    &\leq K_3\int_0^\epsilon e^{-xt}t^{\alpha+\beta(N+1)}\,dt\\
    &\leq K_3\int_0^\infty e^{-xt}t^{\alpha+\beta(N+1)}\,dt\\
    &=K_3\frac{\Gamma(\alpha+\beta(N+1)+1)}
    {x^{\alpha+\beta(N+1)+1}}.
\end{align*}
Since $\beta>0$,
\[
    R_1=o\!\left(x^{-\alpha-\beta N-1}\right)
    \qquad\text{as }x\to\infty.
\]
This is true for every $N$, so
\[
    F_1(x)\sim\sum_{n=0}^\infty
    a_n\frac{\Gamma(\alpha+\beta n+1)}{x^{\alpha+\beta n+1}}.
\]
Since $F_2$ is smaller than every inverse power of $x$, the same expansion
holds for $F=F_1+F_2$:
\[
    F(x)\sim\sum_{n=0}^\infty
    a_n\frac{\Gamma(\alpha+\beta n+1)}{x^{\alpha+\beta n+1}}
    \qquad\text{as }x\to\infty.
\]

\subsection*{3.5 Laplace integrals and Laplace's method}

Laplace integrals are of the form
\[
    I(x)=\int_a^b f(t)e^{x\phi(t)}\,dt,
\]
where $b>a$, and $f$ and $\phi$ are sufficiently nice functions.
What happens to $I$ as $x\to\infty$?
Watson's lemma deals with the special case $\phi(t)=-t$, $a=0$, $b=T$.

In general, we have Laplace's \emph{principle of localisation}:
the asymptotic expansion of $I$ as $x\to\infty$ is determined by
contributions from arbitrarily small regions around the maxima of $\phi$.
Figure~\ref{fig:am-lecture5-localisation} illustrates this principle.
\begin{figure}[htbp]
    \centering
    \pdfximage width 0.88\linewidth {figures/asymptotic-lecture5-localisation.jpg}
    \pdfrefximage\pdflastximage
    \caption{The three blackboard sketches of localisation. Left: Watson's
    case, where the integral is dominated by an endpoint. Upper right:
    a phase $\phi$ with an interior maximum at $t=c$. Lower right:
    the integrand becomes concentrated in a tiny region around $c$.
    Source: lecture blackboard photograph.}
    \label{fig:am-lecture5-localisation}
\end{figure}

\noindent\textbf{(a) Suppose $\phi$ is monotonic on $[a,b]$.}
If $\phi'(t)>0$ and $f(b)\neq0$, let
\[
    u=\phi(t),\qquad t=\phi^{-1}(u).
\]
Then
\[
    I(x)=\int_{\phi(a)}^{\phi(b)}
    \frac{f(t)}{\phi'(t)}e^{xu}\,du,
    \qquad t=\phi^{-1}(u).
\]
Comparing with Example 2 from Section 3.2 gives
\[
    I(x)\sim\frac{f(b)e^{x\phi(b)}}{x\phi'(b)}
    \qquad\text{as }x\to\infty.
\]
If $\phi'(t)<0$ and $f(a)\neq0$, then similarly
\[
    I(x)\sim\frac{f(a)e^{x\phi(a)}}{x|\phi'(a)|}
    \qquad\text{as }x\to\infty.
\]

\noindent\textbf{(b) Suppose $\phi(t)$ has one interior maximum at $t=c$.}
Assume
\[
    \phi'(c)=0,\qquad \phi''(c)<0,\qquad f(c)\neq0.
\]
For small, positive $\epsilon$, write
\[
    \Omega_\epsilon=[a,c-\epsilon]\cup[c+\epsilon,b]
\]
and
\begin{align*}
    I(x)&=\int_{c-\epsilon}^{c+\epsilon}f(t)e^{x\phi(t)}\,dt
    +\int_{\Omega_\epsilon}f(t)e^{x\phi(t)}\,dt\\
    &\equiv I_1(x)+I_2(x).
\end{align*}
The small region around $c$ is the region responsible for the concentration
shown in Figure~\ref{fig:am-lecture5-localisation}.

First, bound $I_2$. In the monotonic regions on either side of $c$,
\begin{align*}
    I_2(x)&=\int_{c+\epsilon}^b f(t)e^{x\phi(t)}\,dt
    +\int_a^{c-\epsilon}f(t)e^{x\phi(t)}\,dt\\
    &=O\!\left(\frac{e^{x\phi(c+\epsilon)}}{x}\right)
    +O\!\left(\frac{e^{x\phi(c-\epsilon)}}{x}\right).
\end{align*}
Here $\phi$ is monotonically decreasing on the right and monotonically
increasing on the left. But
\[
    \phi(c+\epsilon)<\phi(c),\qquad
    \phi(c-\epsilon)<\phi(c),
\]
so
\[
    I_2(x)=o\!\left(\frac{e^{x\phi(c)}}{x}\right)
    \qquad\text{as }x\to\infty.
\]

Now consider $I_1$. Expand $f$ and $\phi$ around $t=c$, assuming they are
sufficiently differentiable:
\begin{align*}
    \phi(t)&=\phi(c)+\underbrace{\phi'(c)(t-c)}_{=0}
    +\tfrac12\phi''(c)(t-c)^2+O((t-c)^3),\\
    f(t)&=f(c)+O(t-c)
    \qquad\text{as }t\to c.
\end{align*}
Thus
\begin{align*}
    I_1(x)
    &=\int_{c-\epsilon}^{c+\epsilon}
    [f(c)+\text{H.O.T.}]\,
    \exp\!\left[x\phi(c)+\tfrac12x\phi''(c)(t-c)^2
    +\text{H.O.T.}\right]dt\\
    &=f(c)e^{x\phi(c)}\int_{c-\epsilon}^{c+\epsilon}
    e^{x\phi''(c)(t-c)^2/2}\,[1+\text{H.O.T.}]\,dt.
\end{align*}
The idea is that the higher-order terms (H.O.T.) in powers of $(t-c)$
generate higher-order terms in an asymptotic series in inverse powers
of $x$ (see Chapling, Section 3.3.1, and B\&O, Section 6.4).
For now, we generate the leading-order term:
\[
    I_1(x)\sim f(c)e^{x\phi(c)}
    \int_{c-\epsilon}^{c+\epsilon}e^{x\phi''(c)(t-c)^2/2}\,dt
    \qquad\text{as }x\to\infty.
\]

\begin{center}
    \textbf{End of lecture 5}
\end{center}

\end{document}

